\scanpage{5}
\deleted{$\tau_0\tau_2=\tau_2\tau_0$} \unsure{i.e.}

Donnée (en plus de $X$ fibré en droites projectives) \inserted{d'une Alg. de quatern. $M$} $\Rightarrow C$ assoc. proj.

\diagram{\resizebox{\linewidth}{!}{\begin{tikzcd}[ampersand replacement=\&,row sep=small,column sep=tiny]
{\deleted{\unsure{\exists}}\ \underbrace{\tau_0\ \tau_1\ \tau_2}}\arrow[d,Leftrightarrow]\&{\in\operatorname{Aut}(X)\quad\text{réflexions }(\text{i.e. }\tau_i^2=\mathrm{id},\ \tau_i\ne1\text{ chaque fibre})}\\
{\overbrace{h_0,h_1,h_2}}\arrow[d,Rightarrow]\&{\in\Gamma(C\setminus X)}\\
{L_0,L_1,L_2}\&{\text{orth. (droites des pts fixes des }\tau_0^C,\tau_1^C,\tau_2^C\text{)}}
\end{tikzcd}}}
\margin{\uline{NB} les $\tau_i$ se verraient de façon canonique \deleted{\unsure{en des}} à un signe près en des \deleted{\unsure{élts de}}\inserted{réflexions} \deleted{\unsure{$G$}}$M^*$}

\uline{Hyp}
\[
\left\{\begin{aligned}
&A)\quad\tau_0\tau_2=\tau_2\tau_0.\quad\text{i.e. } A')\ h_0\perp h_2\\
&B)\quad\tau_0\tau_1\ne\tau_1\tau_0,\quad\tau_1\tau_2\ne\tau_2\tau_1\quad\text{sur chaque fibre}\\
&\qquad\text{ou }\beta_{\unsure{3|1|f}}^2\ne1,\ \beta_{\unsure{5|0|s}}^2\ne1\\
&\qquad\text{i.e. } B')\ h_0\not\perp h_1,\quad h_1\not\perp h_2\quad+\quad\text{---}\\
&\qquad(\text{ou }h_0\notin L_1,\ h_1\notin L_2\quad\text{---})
\end{aligned}\right.
\]
% slanted
\margin{ceci, plus $\tau_i^2=\mathrm{id}$ $(i=0,1,\unsure{2|3})$ impliquent que les $\tau_i$ sont des réflexions (ou $\tau_i^2\ne1$ sur chaque fibre)}
% slanted
\margin{ceci implique aussi, déjà, que $h_i\notin X$ sur chaque fibre $(i=0,1,2)$}

\deleted{C) Les $\tau_0,\tau_1,\tau_2$ \unsure{déf.} \unsure{droites} i.e. pts fixes des $\tau_i$ distincts}

\deleted{C) Les $h_0,h_1,h_2$ engendrent le}

\deleted{C') $h_i\notin\Gamma X$ $(i=0,1,2)$ sur tte fibre}

C) $\{h_0,h_1,h_2\}$ engendre tout $C$ en chaque fibre

% slanted
\margin{comment l'exprimer en termes des opérations $\tau_i$ sur \uline{$X$} ?}
($\Longleftrightarrow H_0,H_1,H_2$ indép. en chaque fibre)

D) Considérons la droite \inserted{$H_1\cap H_2$} de $\check{\mathbb P}(M)\supset C\simeq\check{\mathbb P}(\unsure{\gamma M|\mathscr V M})$ formée des pts orth. à $h_1,\deleted{\unsure{h_0}},h_2$ i.e. fixes par $\tau_1,\tau_2$. \deleted{On voit qu'elle ne sont pas tang} \inserted{Les conditions précédentes impliquent qu'en tte fibre,} elle est non contenue dans la quadrique can. de $\check{\mathbb P}(M)$ (définie par le déterminant).

On exige qu'elles \unsure{ne} \unsure{sont|soit} \unsure{pas} tangentes \inserted{en tte fibre}, ce qui signifie aussi que $\beta\deleted{\unsure{=|\ne}}2$ \unsure{sur|inv.}, ou encore que a) Aux pts \inserted{de $S$} de car. $p>0$, \inserted{impaire,} on veut que l'ordre de $\tau_1\tau_2$ soit $\ne p$ (celà suffit \inserted{sur les pts de car. $>0$ \deleted{sont \unsure{d'ordre} \unsure{puisque} \unsure{contr.}} \inserted{\unsure{tous deux|des deux|deux} \unsure{supp.|pts}} \deleted{\unsure{Compte}} \deleted{il faut \unsure{faire}}} car si $\tau_1\tau_2$ d'ordre fini -- sinon on \unsure{cont|contre} aux pts de car. rés. nulle.
% slanted
\margin{il suffit m\textsuperscript{m} pour tte composante irréductible de $S$ contenant un point de car. $p>0$ !}
% archive sheet number: 4
